\documentclass{article}
\usepackage{graphicx}
\usepackage{listings}
\usepackage{color}
\usepackage{amsmath}
\usepackage{hyperref}
\usepackage{tikz}
\usetikzlibrary{positioning}
\usepackage{tikz-uml}
\usepackage[margin=1in]{geometry}

\title{Project Calculator\\
\large Arbitrary precision arithmetic Library}

\date{May 2025}

\begin{document}

\maketitle
\section*{Directory Structure}
\begin{verbatim}
project_java/
├── build.xml
├── src/
│   ├── arbitraryarithmetic/
│   │    ├── AInteger.java
│   │    ├── AFloat.java
│   └── MyInfArith.java
├── dist/  aarithmetic.jar
├── run_arith.py
├── report.pdf
├── README.md
\end{verbatim}

\section*{Introduction}
This report describes the design section, which describes the design decisions that I made and also include the class diagrams in UML format.


\section*{Team Details}
\begin{itemize}
  \item Name: Petei Divya Poorna Chand
  \item Roll Number: CS24BTECH11048
\end{itemize}

\section*{Design}
The central idea of this project is to perform arithmetic operations on Integers and Floating-Point numbers without using libraries Bignumbers or any other built-in types
\begin{itemize}
    \item \texttt{AInteger}:Handles arbitrarily large integers, including sign management.
    \item \texttt{AFloat}:  Supports floating-point arithmetic by separating integer and decimal parts and using \texttt{AInteger} for computations.
\end{itemize}
\section*{Design Decisions}

This section explains the key design decisions made during the development of the Arbitrary Precision Arithmetic Library.

\subsection*{1. Object-Oriented Approach}
We adopted an object-oriented programming (OOP) paradigm to structure the library. Two main classes, \texttt{AInteger} and \texttt{AFloat}, were created to encapsulate the behavior and data of arbitrary-precision integers and floating-point numbers, respectively.

\subsection*{2. String-Based Internal Representation}
To support numbers of arbitrary size and precision, we chose to store numeric values as \texttt{String} objects. This allowed us to overcome the limitations of built-in Java data types like \texttt{int}, \texttt{long}, or \texttt{double}, and manually implement arithmetic logic character by character.

\subsection*{3. Manual Arithmetic Implementation}
All arithmetic operations (+, –, $\times$, $\div$) were implemented manually to simulate how operations are performed on paper:
\begin{itemize}
    \item Addition and subtraction follow digit-wise traversal with carry/borrow handling.
    \item Multiplication uses column-based manual multiplication.
    \item Division follows the long division method, including repeated subtraction and decimal handling.
\end{itemize}
This design ensures high precision and full control over the operations.

\subsection*{4. Separation of Integer and Floating-Point Logic}
The decision to use separate classes for \texttt{AInteger} and \texttt{AFloat} ensures clarity in handling integer-only logic versus floating-point arithmetic:
\begin{itemize}
    \item \texttt{AInteger} handles whole numbers and is optimized for base operations.
    \item \texttt{AFloat} wraps two \texttt{AInteger} objects (for integer and fractional parts) and manages decimal point alignment.
\end{itemize}

\subsection*{5. Command Line Interface (CLI)}
We provided a CLI through the \texttt{MyInfArith} class. This allows users to perform operations directly from the terminal, taking two number strings and an operator as input and printing the result to the console.

\subsection*{6. Ant-Based Build System}
To streamline compilation and packaging, we used Apache Ant to define a build process via \texttt{build.xml}. This makes the project easy to build and test with a single command.
\texttt{aarithmetic.jar} file.

\subsection*{7. Decoupling Testing and Execution}
Testing and execution are decoupled using an external Python script. The script builds the project and executes the JAR file with user inputs, mimicking test case runs in a controlled environment.

\subsection*{8. Precision Control}
To limit performance while preserving accuracy, we designed the system to support up to 30 decimal digits of precision for floating-point operations. This level of precision is adequate for most use cases and balances speed with accuracy.



\section*{implementation}
\subsection*{AInteger class}
The \texttt{AInteger} class handles arbitrary-precision integers. All numbers are represented as strings to avoid Java's built-in limits. Operations supported include:
\begin{itemize}
    \item \texttt{add(AInteger other)} – Adds two AIntegers using string-based digit-wise addition,Handling carry and sign.
    \item \texttt{subtract(AInteger other)} – Performs subtraction, handling borrowing and sign.
    \item \texttt{multiply(AInteger other)} – Repeated addition of one number using AInteger.add(Other).
    \item \texttt{divide(AInteger other)} – Repeated subtraction of one number from other using AInteger.subtract(other) and counts repetitions
\end{itemize}

\subsection*{AFloat Class}
The \texttt{AFloat} class handles floating-point numbers using \texttt{AInteger} for the integer parts. It supports:
\begin{itemize}
    \item stores index of decimal point and makes floating point numbers into integers by removing decimal points
    \item uses AInteger for operations and add the decimal point at required place based on operation
    \item Manual multiplication and division with precision up to 30 decimal digits.
    \item sign handling
\end{itemize}

\section*{Command-Line Interface:MyInfArith Class}
This class acts as the main interface for users. It parses command-line arguments and implements the appropriate arithmetic operations.
\subsection*{Functionality}
\begin{itemize}
    \item Accepts input in the format: \texttt{java MyInfArith type operation num1 num2}
    \item Determines whether the operands are integers or floating-point based on the presence of a decimal point.
    \item Calls the corresponding method from \texttt{AInteger} or \texttt{AFloat} for the requested operation.
    \item Outputs the result.
\end{itemize}
\section*{Creating a jar file using ant}
I used Apache Ant, a Java-based build automation tool, to compile our source code and package it into a JAR file named \texttt{aarithmetic.jar}. This approach allows for reproducible builds and eliminates manual compilation and packaging.
\subsection*{Buildfile : build.xml}
    The build configuration is defined in a file named \texttt{build.xml}. This file specifies the project structure, compilation target, and JAR creation logic.
\subsection*{How to Build}
To build and generate the JAR file, simply run:
\begin{verbatim}
ant jar
\end{verbatim}
in terminal This will:
\begin{itemize}
    \item Compile all Java files under \texttt{src/arbitraryarithmetic/}
    \item Place compiled \texttt{.class} files into \texttt{target/}
    \item Create a JAR file \texttt{aarithmetic.jar} with \texttt{MyInfArith} as the main class
\end{itemize}
\subsection*{Advantages}
Using Ant simplifies the build process by:
\begin{itemize}
    \item Automating compilation and packaging
    \item Ensuring consistent builds
    \item Making the project more portable and easier to test
\end{itemize}

\section*{python code to run the main program}
To reduce difficulty in the the development and execution process, we created a Python script that automates compiling the Java source files and running the resulting executable JAR file. This script significantly reduces manual effort and ensures a consistent build environment for testing and demonstration purposes.
\subsection*{Functionality}
The Python script performs the following actions:
\begin{enumerate}
    \item It compiles all the Java files in the \texttt{src/arbitraryarithmetic/} directory using the \texttt{javac} compiler.
    \item It packages the compiled \texttt{.class} files into a JAR named \texttt{aarithmetic.jar}, setting \texttt{MyInfArith} as the main class.
    \item It executes the JAR file with user-provided arguments (  \texttt{type},\texttt{operator},\texttt{num1}, \texttt{num2}) and prints the output.
\end{enumerate}
\subsection*{Usage}
\textbf{Command-line format:}
\begin{verbatim}
python3 run_arithmetic.py <type> <operator> <num1> <num2>
\end{verbatim}

\section*{Class Diagram (UML)}

\begin{center}
\begin{tikzpicture}
  \umlclass[x=-10, y=0]{AInteger}{
    - value : String
  }{
    + AInteger() \\
    + AInteger(value : String) \\
    + copy(other : AInteger) : AInteger \\
    + parse(s : String) : AFloat\\
    + toString() : String \\
    + add(other : AInteger) : AInteger \\
    + subtract(other : AInteger) : AInteger \\
    + multiply(other : AInteger) : AInteger \\
    + divide(other : AInteger) : AInteger \\
    + compare(String a,String b) : int
   \\
  }

  \umlclass[x=0, y=0]{AFloat}{
    - value : String
  }{
    + AFloat() \\
    + AFloat(value : String) \\
    + copy(other : AFloat) : AFloat \\
    + parse(s : String) : AFloat\\
    + toString() : String \\
    + add(other : AFloat) : AFloat \\
    + subtract(other : AFloat) : AFloat \\
    + multiply(other : AFloat) : AFloat \\
    + divide(other : AFloat) : AFloat \\
    + compare(String a,String b) : int
    }
    \uml{AFloat}{AInteger}
\end{tikzpicture}
\end{center}
\subsection*{Limitations}

\begin{enumerate}
    \item \textbf{Manual Long Division is Time-Consuming:} \\
    The \texttt{divide} methods in both \texttt{AInteger} and \texttt{AFloat} use a manual, long-division-style approach implemented through string manipulation. As a result, the execution time increases significantly for large input values.

    \item \textbf{Fixed Precision of 30 Decimal Digits:} \\
    The \texttt{AFloat} class supports up to 30 digits after the decimal point. Any computation requiring higher precision will be truncated or rounded beyond this limit.
\end{enumerate}


\end{document}
